%==========================================================
% ANALYSIS: Part 1 table
%==========================================================
\paragraph{} In part 1 we created created a three phase transformer bank in the $\Delta$-wye configuration. The LabVolt 8348
transformer banks in the lab are rated at 250VA with a turns ratio of 1 if the center tap is not in use. When we compare values
between the measurements in table 1 we see a pattern, but also some differences. Our measurements and calculations are off by
a half in some places. This is probably be explanable. But I'm sleepy.

\begin{table}[htbp]
	\centering
	\caption{Part 1: $\Delta$-Wye bank, measured vs. calculated (rated values)}
	\label{tab:part1}
	\begin{tabular}{lllrrr}
		\toprule
		\textbf{LabVolt Label} & \textbf{Symbol} & \textbf{Unit} & \textbf{Measured} & \textbf{Calculated} & \textbf{Error (\%)} \\
		\midrule
		V Bank Primary         & $V_{Bank,P}$    & V             & 207.3             & 208.0               & $0.34$              \\
		V Phase Primary        & $V_{Phase,P}$   & V             & 207.1             & 208.0               & $0.43$              \\
		V Bank Secondary       & $V_{Bank,S}$    & V             & 200.3             & 207.85              & $3.63$              \\
		V Phase Secondary      & $V_{Phase,S}$   & V             & 116.0             & 120.0               & $3.33$              \\
		I Bank Primary         & $I_{Bank,P}$    & A             & 1.013             & 1.20                & $18.4$              \\
		I Phase Primary        & $I_{Phase,P}$   & A             & 0.585             & 0.693               & $18$                \\
		I Bank Secondary       & $I_{Bank,S}$    & A             & 0.973             & 1.20                & $23$                \\
		I Phase Secondary      & $I_{Phase,S}$   & A             & 0.742             & 0.693               & $-6.6$              \\
		PQS1 (E1, I1) 3$\sim$  & $S_{3P}$        & VA            & 428               & 432                 & $0.01$              \\
		\bottomrule
	\end{tabular}
\end{table}


%==========================================================
% ANALYSIS: Part 2 tables
%==========================================================

\paragraph{} In Part 2 of this lab, we configured a single-phase transformer to replicate a residential distribution system.
The primary side was supplied with 208 V from a three-phase source, while the secondary side supplied two 120 V legs that were 180$^\circ$
out of phase, allowing us to connect both 120 V and 240 V loads.

Table 2 shows the measured and calculated values with all three loads connected. We observed a decrease in secondary voltages of approximately 4.5\%
to 5.6\%, likely due to the internal impedance of the transformer. The two 120 V loads drew nearly identical currents, resulting in a neutral current
of only 0.009 A rather than the calculated 0 A. This difference could be due to slight variations in the loads or measuring equipment.
\clearpage
\begin{table}[htbp]
	\centering
	\caption{Part 2: Full load, measured (average of 3 readings) vs. calculated}
	\label{tab:part2full}
	\begin{tabular}{lllrrr}
		\toprule
		\textbf{LabVolt Label} & \textbf{Symbol} & \textbf{Unit} & \textbf{Measured} & \textbf{Calculated} & \textbf{Error (\%)} \\
		\midrule
		Line Voltage           & $V_{Line}$      & V             & 207.5             & 208.0               & $-0.24$             \\
		H1 Voltage             & $V_{H1}$        & V             & 114.6             & 120.0               & $-4.50$             \\
		H2 Voltage             & $V_{H2}$        & V             & 113.3             & 120.0               & $-5.58$             \\
		240 Voltage            & $V_{240}$       & V             & 227.2             & 240.0               & $-5.33$             \\
		Line Current           & $I_{Line}$      & A             & 0.250             & 0.265               & $-5.7$              \\
		H1 Current             & $I_{H1}$        & A             & 0.090             & 0.098               & $-8.2$              \\
		H2 Current             & $I_{H2}$        & A             & 0.089             & 0.098               & $-9.2$              \\
		Neutral Current        & $I_{N}$         & A             & 0.009             & 0                   & --                  \\
		Line Power             & $P_{Line}$      & W             & 51.80             & 47.08               & $+10.0$             \\
		H1 Power               & $P_{H1}$        & W             & 10.34             & 11.54               & $-10.4$             \\
		H2 Power               & $P_{H2}$        & W             & 10.11             & 11.54               & $-12.4$             \\
		\bottomrule
	\end{tabular}
\end{table}

\subsection{Full Load}
\paragraph{}Table 2 shows the measured and calculated values with all three loads connected. We observed a decrease in secondary voltages of approximately 4.5% to 5.6%, likely due to the internal impedance of the transformer. The two 120 V loads drew nearly identical currents, resulting in a neutral current of only 0.009 A rather than the calculated 0 A. This difference could be due to slight variations in the loads or measuring equipment.

\begin{table}[htbp]
	\centering
	\caption{Part 2: No load, measured (average of 3 readings) vs. calculated}
	\label{tab:part2noload}
	\begin{tabular}{lllrrr}
		\toprule
		\textbf{LabVolt Label} & \textbf{Symbol} & \textbf{Unit} & \textbf{Measured} & \textbf{Calculated} & \textbf{Error (\%)} \\
		\midrule
		Line Voltage           & $V_{Line}$      & V             & 207.8             & 208.0               & $-0.10$             \\
		H1 Voltage             & $V_{H1}$        & V             & 120.1             & 120.0               & $+0.08$             \\
		H2 Voltage             & $V_{H2}$        & V             & 119.5             & 120.0               & $-0.42$             \\
		240 Voltage            & $V_{240}$       & V             & 239.1             & 240.0               & $-0.38$             \\
		Line Current           & $I_{Line}$      & A             & 0.017             & 0                   & --                  \\
		H1 Current             & $I_{H1}$        & A             & 0.008             & 0                   & --                  \\
		H2 Current             & $I_{H2}$        & A             & 0.008             & 0                   & --                  \\
		Neutral Current        & $I_{N}$         & A             & 0.009             & 0                   & --                  \\
		Line Power             & $P_{Line}$      & W             & 3.43              & 0                   & --                  \\
		H1 Power               & $P_{H1}$        & W             & 0.940             & 0                   & --                  \\
		H2 Power               & $P_{H2}$        & W             & 0.907             & 0                   & --                  \\
		\bottomrule
	\end{tabular}
\end{table}
%
\subsection{No Load}
\paragraph{}Table 3 shows the results with all loads disconnected. Without loading on the secondary side, the voltages increased to within 0.5% of their calculated values. A small amount of current and power was still measured on the primary side due to the magnetization of the core and transformer losses. Small measurements on the secondary side were likely due to limitations of the measuring equipment. Percentage errors could not be calculated where the expected values were 0.

\begin{table}[htbp]
	\centering
	\caption{Part 2: Imbalanced load (H1 load open), measured vs. calculated}
	\label{tab:part2imb}
	\begin{tabular}{lllrrr}
		\toprule
		\textbf{LabVolt Label} & \textbf{Symbol} & \textbf{Unit} & \textbf{Measured} & \textbf{Calculated} & \textbf{Error (\%)} \\
		\midrule
		Line Voltage           & $V_{Line}$      & V             & 207.2             & 208.0               & $-0.38$             \\
		H1 Voltage             & $V_{H1}$        & V             & 116.4             & 120.0               & $-3.00$             \\
		H2 Voltage             & $V_{H2}$        & V             & 114.5             & 120.0               & $-4.58$             \\
		240 Voltage            & $V_{240}$       & V             & 229.8             & 240.0               & $-4.25$             \\
		Line Current           & $I_{Line}$      & A             & 0.205             & 0.213               & $-3.6$              \\
		H1 Current             & $I_{H1}$        & A             & 0.008             & 0                   & --                  \\
		H2 Current             & $I_{H2}$        & A             & 0.090             & 0.098               & $-8.2$              \\
		Neutral Current        & $I_{N}$         & A             & 0.088             & 0.098               & $-10.3$             \\
		Line Power             & $P_{Line}$      & W             & 42.42             & 35.54               & $+19.4$             \\
		H1 Power               & $P_{H1}$        & W             & 0.92              & 0                   & --                  \\
		H2 Power               & $P_{H2}$        & W             & 10.32             & 11.54               & $-10.6$             \\
		\bottomrule
	\end{tabular}
\end{table}

\subsection{Imbalanced Load}
\paragraph{}Table 4 shows the results after disconnecting the 120 V load on H1. With less loading, we observed a slight increase in secondary voltages and a decrease in primary current and power. The most noticeable difference was the increase in neutral current from 0.009 A to 0.088 A. This was expected as the current from the remaining 120 V load was no longer being canceled out by the opposing load.
Through these tests, we were able to observe the effects of loading on the secondary voltage of a transformer and the importance of balancing loads within a residential distribution system to minimize neutral current.

\clearpage
